% ID: SEL-D0077-OB-E02
% Tanda S077.
\textbf{a)} El dato $f(2)=7$ implica $2a+b=7$. Para que haya continuidad en $1$, el límite derecho $a+b$ debe coincidir con $f(1)=1^2+4=5$. Restando las ecuaciones obtenemos $\boxed{a=2,\ b=3}$.

\textbf{b)} En un entorno de $-1$ se utiliza el tramo $x^2+4$, cuya derivada es $2x$. El punto de tangencia es $(-1,5)$ y la pendiente es $-2$. Aplicamos la ecuación punto-pendiente:
\[
 y-5=-2(x+1),\qquad\boxed{y=-2x+3}.
\]
